\section{Supplementary material}

\subsection{Agreement of the grid and the paired runs}
The full grid and the paired W3/G128 exports were produced by separate evaluation runs of the same checkpoints.
On the \V{consistContexts} contexts present in both, baseline top-1 differs by at most \V{consistFp} points and W3/G128 top-1 by at most \V{consistQ} points.

\subsection{Matching bin width}
With bin widths 0.025, 0.05 and 0.1, the reduction of the retention gap ranges from \V{reductionMinAll}\% to \V{reductionMaxAll}\% over the nine settings.

\subsection{Correlates of the difficulty score}
\Cref{tab:properties} lists the Spearman correlation of every measured property with the difficulty score.

\begin{table}[h]
\centering
\caption{Spearman correlation of each property with the difficulty score over the {\V{gridDatasets}} complete datasets. FP: baseline top-1. Sep.: separability. H: H-score. Norm: feature norm ratio P99/P50. All are measured on the baseline except the final-layer CKA between baseline and W3/G128 features.}
\label{tab:properties}
\setlength{\tabcolsep}{1.5pt}
\footnotesize
\begin{tabular}{llrrrrrrr}
\toprule
Backbone & Regime & $q_{10}$ & FP & $C$ & Sep. & H & Norm & CKA \\
\midrule
ViT-B/16 & FT & $-$0.98 & $-$0.97 & 0.52 & $-$0.84 & $-$0.32 & 0.42 & $-$0.28 \\
ViT-B/16 & FZ & $-$0.92 & $-$0.90 & 0.37 & $-$0.83 & $-$0.24 & $-$0.90 & $-$0.55 \\
ViT-B/16 & CAL & $-$0.83 & $-$0.83 & 0.48 & 0.06 & $-$0.03 & 0.41 & $-$0.07 \\
Swin-B & FT & $-$0.86 & $-$0.78 & 0.55 & $-$0.73 & $-$0.58 & $-$0.37 & $-$0.60 \\
Swin-B & FZ & $-$0.97 & $-$0.98 & 0.44 & $-$0.81 & $-$0.15 & $-$0.80 & $-$0.33 \\
Swin-B & CAL & $-$0.63 & $-$0.59 & 0.44 & 0.07 & $-$0.51 & $-$0.11 & $-$0.18 \\
BEiTv2-B & FT & $-$0.48 & $-$0.40 & 0.07 & $-$0.39 & $-$0.21 & $-$0.16 & $-$0.36 \\
BEiTv2-B & FZ & $-$0.92 & $-$0.95 & 0.36 & $-$0.88 & $-$0.43 & 0.70 & $-$0.11 \\
BEiTv2-B & CAL & $-$0.83 & $-$0.74 & 0.16 & 0.30 & $-$0.06 & $-$0.08 & 0.28 \\
\bottomrule
\end{tabular}
\end{table}


\subsection{Local structure of the features}
\Cref{fig:tsne} shows a joint t-SNE~\cite{vandermaaten2008tsne} of baseline and W3/G128 final features for one context.
In this context, the clusters of the selected classes remain separate after quantization.

\begin{figure}[h]
  \centering
  \includegraphics[width=\linewidth]{figures/tsne_cub_vit_ft.pdf}
  \caption{Joint t-SNE of the final features of ViT-B/16 FT on CUB, for \V{tsneClasses} of its \V{cubClasses} classes and all their test images, before and after W3/G128 quantization. Colors are classes.}
  \label{fig:tsne}
\end{figure}
